\SourceLeaf{21}

Die Observatio kann aus den Criterien leicht abgeleitet werden. Z.\,B. wurde für den kubischen Rest 3
soeben die Bedingung $N\equiv 0\;(\mathrm{mod.}\ 3)$ gewonnen und unter [2] ebenso angegeben. Sind nun $M$ und $N$
gerade, so setze man
\[ M=2x,\qquad N=6y, \]
sodass aus $4n=M^2+27N^2$ die Darstellbarkeit von $n$ durch die quadratische Form mit der Determinante $-243$
\[ n=x^2+243y^2 \]
mit ganzzahligen $x, y$ folgt. Umgekehrt folgt aus der Möglichkeit einer solchen Darstellung von $n$
die Gültigkeit der Kongruenz $N\equiv 0\;(\mathrm{mod.}\ 3)$. Sind $M$ und $N$ ungerade, so wähle man das Vorzeichen richtig und setze
\[ M\pm\dfrac{N}{3}=4x,\qquad \mp\dfrac{N}{3}=y, \]
woraus man die Darstellung
\[ n=4x^2+2xy+61y^2. \]
gewinnt. Aus dieser Formel würde umgekehrt $N\equiv 0\;(\mathrm{mod.}\ 3)$ folgen, sodass 3 genau dann kubischer Rest
modulo $n$ ist, wenn $n$ durch eine der Formen $(1,0,243)$, $(4,1,61)$ darstellbar ist.

Herr R. \textsc{Dedekind}, dem ich mehrere für mich sehr förderliche Unterhaltungen über die in
Rede stehenden Gegenstände verdanke, hat mich darauf aufmerksam gemacht, dass für den Fall des kubischen
Nichtrests 3 modulo $n$ die Angaben von \textsc{Gauss} nicht ganz vollständig sind. Es gibt insgesamt neun Klassen ursprünglicher Formen erster Art mit der Determinante $D=-243$, deren reduzierte Formen die folgenden sind
\[ (1,0,243),\quad (4,\pm1,61),\quad (13,\pm2,19),\quad (7,\pm3,36),\quad (9,\pm3,28). \]
In [1] fehlen die Formen $(13,\pm2,19)$. Das Theorem muss demnach so lauten, dass 3 genau
dann kubischer Nichtrest modulo $n$ ist, wenn $n$ eine Darstellung in einer der folgenden Formen gestattet
\[ n=13xx+4xy+19yy, \]
\[ n=7xx+6xy+36yy, \]
\[ n=9xx+6xy+28yy. \]
Diese Darstellungen erhält man stets, indem man bei den folgenden Transformationen das Vorzeichen richtig wählt
\[ \text{I. Wenn } M\equiv N\equiv 1\;(\mathrm{mod.}\ 2),\quad M\pm 7N\equiv 0\;(\mathrm{mod.}\ 12), \]
\[ M=5x+7y,\qquad \pm N=x-y. \]
\[ \text{II. Wenn } M\equiv N\equiv 1\;(\mathrm{mod.}\ 2),\quad M\pm 7N\equiv 6\;(\mathrm{mod.}\ 12), \]
\[ M=x+6y,\qquad \mp N=x. \]
\[ \text{III. Wenn } M\equiv N\equiv 0\;(\mathrm{mod.}\ 2),\quad M\pm N\equiv 0\;(\mathrm{mod.}\ 3), \]
\[ M=6x+2y,\qquad \mp N=2y. \]
\par\noindent\hspace*{\parindent}Die Angaben zur Zahl 5 unter 1 folgen in ähnlicher Weise aus dem Restkriterium
$MN\equiv 0\;(\mathrm{mod.}\ 5)$.

Im Übrigen sei noch bemerkt, dass sich im Original der Notiz [1] eine bis 223 fortgesetzte Tabelle derjenigen Primzahlen $p=3n+1$ findet, für die die Zahlen 2, 3, 5, 7 kubische Reste sind.

\hfill \textsc{Fricke.}

\endgroup
